\documentclass[UTF8,fontset=none,11pt,a4paper]{ctexart}
\usepackage{ipara-handbook}
\handbooksetup{NET-I01 一个网络请求的一生}{408 · 计算机网络 Integration}{请求过程 · 逐跳交付 · 状态 · 边界}
\begin{document}
\handbooktitle{一个网络请求的一生}{从输入 URL 到收到网页内容}{408 · NET-I01}

\begin{corebox}{这本综合册要解决什么问题}
浏览器里只输入了一个 URL，但网络不会因此机械执行一遍“DHCP--DNS--ARP--TCP--HTTP”。主机可能已经有网络配置，DNS、ARP、TCP 或 HTTP 缓存也可能已经存在；不同链路还会使用不同的帧格式。

本册要回答的是：\textbf{在当前已有状态和当前网络拓扑下，为了完成这次访问，下一步真正需要做什么；每产生一条数据报，它又怎样一跳一跳到达目标。}
\end{corebox}

\begin{methodbox}{本册负责组合，不重复讲协议细节}
NET01--NET08 分别负责时延、单跳交付、可靠传输、IP 转发、路由、TCP/UDP、拥塞控制和应用层协议；NET-B01--B06 负责这些专题之间的稳定接口。

本册只把它们接成一个完整过程：什么状态可以直接复用，什么状态必须先建立；一条数据报经过主机、AP、交换机、路由器和不同链路时，哪些字段保持、哪些字段改变；失败时第一处断点在哪里。
\end{methodbox}

\section{先建立正确的总图：请求推进和逐跳交付是两个不同过程}

一个容易产生的错误画法是：
\[
\text{DHCP}\to\text{DNS}\to\text{ARP}\to\text{TCP}\to\text{HTTP}.
\]
这条线只能当知识点清单，不能当真实执行顺序。原因很直接：DNS 查询本身就要先被送到 DNS 服务器；TCP 的 SYN、ACK 和 HTTP 数据也都要经过 IP 转发和当前链路交付。ARP 只在需要解析当前 Ethernet 下一跳时出现，Native PPP 链路根本不需要 Ethernet MAC。

因此把全过程分成两个独立又协同的维度（如图~\ref{fig:request-and-datagram} 所示）：

\begin{enumerate}
  \item \textbf{请求推进（事务层）：}这次访问还缺什么状态？已有状态直接复用，只补真正缺失的部分。
  \item \textbf{逐跳交付（数据报层）：}一旦已经形成一条普通 IP 数据报，当前节点怎样根据自己的 FIB 选出下一跳，并用当前链路把它送过去？
\end{enumerate}

\begin{figure}[!htbp]
  \centering
  \resizebox{0.92\linewidth}{!}{% ── NET-I01 Canonical Architecture: Transaction & Per-Datagram Execution Model ──
\begin{tikzpicture}[
  scale=1.0,
  >=Stealth,
  font=\small,
  color=themefg,
  text=themefg,
  txnnode/.style={draw=themecurve,fill=themecurve!15!themebg,rounded corners=4pt,minimum width=2.1cm,minimum height=1.05cm,align=center,font=\small\bfseries,line width=1.1pt},
  localnode/.style={draw=themegreen,fill=themegreen!15!themebg,rounded corners=4pt,minimum width=2.1cm,minimum height=1.05cm,align=center,font=\small\bfseries,line width=1.1pt},
  pktnode/.style={draw=themeamber,fill=themeamber!16!themebg,rounded corners=4pt,minimum width=2.1cm,minimum height=1.05cm,align=center,font=\small\bfseries,line width=1.1pt},
  branchnode/.style={draw=themealert,fill=themealert!15!themebg,rounded corners=3pt,minimum width=2.2cm,minimum height=0.88cm,align=center,font=\scriptsize\bfseries,line width=1.0pt},
  flowedge/.style={->,draw=themefg!85,line width=1.1pt},
  invokeline/.style={->,draw=themecurve,line width=1.1pt,dashed},
  loopline/.style={->,draw=themeamber,line width=1.1pt,densely dashed}
]

% ══════════════════════════════════════════════════════════════
% ── 1. 6 个核心列的 X 坐标定义（等距 3.4cm，留出 1.3cm 净间隙）──
% ══════════════════════════════════════════════════════════════
\def\colA{1.4}
\def\colB{4.8}
\def\colC{8.2}
\def\colD{11.6}
\def\colE{15.0}
\def\colF{18.4}

% ══════════════════════════════════════════════════════════════
% ── 2. 背景语义分层底盒 ──
% ══════════════════════════════════════════════════════════════

\begin{pgfonlayer}{background}
  % 事务层 (天青蓝)
  \fill[fill=themecurve!7!themebg,draw=themecurve!50,dashed,line width=1.0pt,rounded corners=8pt]
    (-0.3, 7.8) rectangle (20.1, 14.1);

  % 数据报层 (琥珀金)
  \fill[fill=themeamber!7!themebg,draw=themeamber!50,dashed,line width=1.0pt,rounded corners=8pt]
    (-0.3, 1.8) rectangle (20.1, 7.3);

  % 底部结论栏
  \fill[fill=themecurve!8!themebg,draw=themecurve!50,line width=0.9pt,rounded corners=5pt]
    (-0.3, -0.2) rectangle (20.1, 1.3);
\end{pgfonlayer}

\begin{pgfonlayer}{foreground}
  % 区域大标题（数据报层标题起点偏移至 x = 2.4，杜绝与 x = 1.4 的下沉注入连线相交）
  \node[anchor=west,font=\normalsize\bfseries,color=themecurve] at (0.1, 13.5) {【这次访问还缺什么（事务状态机）】已有状态直接复用，只补真正缺失的步骤};
  \node[anchor=west,font=\normalsize\bfseries,color=themeamber] at (2.4, 6.7) {【已经要发送一条 IP 数据报（逐跳交付流水线）】查当前 FIB，确定下一跳，再完成当前这一跳};
\end{pgfonlayer}

% ══════════════════════════════════════════════════════════════
% ── 3. 节点定义（全图节点声明）──
% ══════════════════════════════════════════════════════════════

% 3.1 事务层主干节点 (y = 10.7)
\node[localnode] (intent) at (\colA, 10.7) {① 用户操作\\\scriptsize 触发网页访问};
\node[localnode] (cache)  at (\colB, 10.7) {② 检查本地缓存\\\scriptsize 能否直接满足?};
\node[txnnode]   (config) at (\colC, 10.7) {③ 本机网络配置\\\scriptsize IP与网关就绪?};
\node[txnnode]   (dns)    at (\colD, 10.7) {④ 确认目标 IP\\\scriptsize 已知目标 IP?};
\node[txnnode]   (tcp)    at (\colE, 10.7) {⑤ 建立传输连接\\\scriptsize TCP 连接就绪?};
\node[localnode] (http)   at (\colF, 10.7) {⑥ 应用层交互\\\scriptsize 发送HTTP并收响应};

% 3.2 事务层分支节点
\node[branchnode,draw=themegreen,fill=themegreen!18!themebg] (done) at (\colB, 12.5) {直接读取本地缓存渲染\\\scriptsize 不再发起网络请求};
\node[branchnode,draw=themeamber,fill=themeamber!18!themebg] (dhcp) at (\colC, 9.1) {触发 DHCP 动态配置\\\scriptsize 获取本机 IP/网关/DNS};

% 3.3 数据报层流水线节点 (y = 4.7)
\node[pktnode] (dst)   at (\colA, 4.7) {① 提取目的 IP\\\scriptsize 确定数据报目标};
\node[pktnode] (fib)   at (\colB, 4.7) {② 查找路由表(FIB)\\\scriptsize 最长前缀匹配(LPM)};
\node[pktnode] (nh)    at (\colC, 4.7) {③ 确定下一跳\\\scriptsize 选定出接口与 IP};
\node[pktnode] (arp)   at (\colD, 4.7) {④ 准备链路交付\\\scriptsize Ethernet 查邻居\\\scriptsize PPP 直接到对端};
\node[pktnode] (frame) at (\colE, 4.7) {⑤ 当前链路封装\\\scriptsize 建立本跳 frame/FCS};
\node[pktnode] (hop)   at (\colF, 4.7) {⑥ 下一节点接收\\\scriptsize 校验并完成当前跳};

% 3.4 目标终点主机交付节点
\node[localnode,draw=themegreen,fill=themegreen!18!themebg,minimum width=2.3cm,minimum height=0.90cm] (deliver) at (\colF, 2.6) {到达目标服务器\\\scriptsize 校验无误剥离首部交付上层};

% ══════════════════════════════════════════════════════════════
% ── 4. 连线与下沉总线调度 ──
% ══════════════════════════════════════════════════════════════

% 4.1 事务层主干推进（短文本紧凑标注，彻底消除字与框重叠）
\draw[flowedge] (intent) -- (cache);
\draw[flowedge] (cache) -- node[above=1pt,font=\tiny\bfseries,color=themealert,align=center]{未命中} (config);
\draw[flowedge] (config) -- node[above=1pt,font=\tiny\bfseries,color=themegreen,align=center]{有效} (dns);
\draw[flowedge] (dns) -- node[above=1pt,font=\tiny\bfseries,color=themegreen,align=center]{已获 IP} (tcp);
\draw[flowedge] (tcp) -- node[above=1pt,font=\tiny\bfseries,color=themegreen,align=center]{已连接} (http);

% 4.2 事务层分支连线
\draw[flowedge,draw=themegreen] (cache) -- node[right=2pt,font=\tiny\bfseries,color=themegreen]{强缓存有效} (done);
\draw[invokeline,draw=themealert] (config) -- node[right=2pt,font=\tiny\bfseries,color=themealert]{未配置/过期} (dhcp);

% 4.3 报文触发下沉与水平总线 (y = 7.7)
\draw[invokeline] (dns.south) -- node[right=1pt,font=\tiny,color=themecurve]{DNS 查询} (\colD, 7.7);
\draw[invokeline] (tcp.south) -- node[right=1pt,font=\tiny,color=themecurve]{TCP 握手} (\colE, 7.7);
\draw[invokeline] (http.south) -- node[right=1pt,font=\tiny,color=themecurve]{HTTP 请求} (\colF, 7.7);

\draw[themecurve,line width=1.2pt,dashed] (\colA, 7.7) -- node[above=2pt,font=\scriptsize\bfseries,color=themecurve]{DNS / TCP / HTTP 产生普通 IP 数据报时，统一进入底层逐跳交付流水线} (\colF, 7.7);

% 4.4 总线注入下层入口（精准连接到 ① 节点顶部）
\draw[flowedge,line width=1.4pt,draw=themecurve] (\colA, 7.7) -- node[left=2pt,font=\tiny\bfseries,color=themecurve]{下沉注入} (dst.north);

% 4.5 数据报层流水线推进
\draw[flowedge] (dst) -- (fib);
\draw[flowedge] (fib) -- (nh);
\draw[flowedge] (nh) -- (arp);
\draw[flowedge] (arp) -- (frame);
\draw[flowedge] (frame) -- (hop);

% 4.6 路由器多跳转发循环（平滑圆角回折线）
\draw[loopline] (hop.south) -- ++(0,-0.8) -| node[pos=0.25,below=2pt,font=\scriptsize\bfseries,color=themeamber] {若为中间路由器：TTL 减 1 $\to$ 重算首部校验和 $\to$ 循环查表转发下一跳} (fib.south);

% 4.7 终点主机交付
\draw[flowedge,draw=themegreen] (hop) -- node[right=2pt,font=\tiny\bfseries,color=themegreen]{已达目标} (deliver);

% ══════════════════════════════════════════════════════════════
% ── 5. 底部核心结论底栏 ──
% ══════════════════════════════════════════════════════════════

\node[anchor=west,font=\normalsize\bfseries,color=themecurve] at (0.1, 0.9) {★ 这张图只表达两个核心关系：};
\node[anchor=west,font=\footnotesize\bfseries,color=themefg] at (0.1, 0.4) {
  1. \textbf{先看已有状态（事务层）：}强缓存、网络配置、DNS 结果、TCP 连接能复用就复用，不机械重跑协议；
};
\node[anchor=west,font=\footnotesize\bfseries,color=themefg] at (0.1, 0.0) {
  2. \textbf{再看每条数据报（逐跳层）：}查当前 FIB、确定下一跳，再按当前链路类型封装完成这一跳；路由器逐跳重封装，二层设备不减 TTL。
};

\end{tikzpicture}
}
  \caption{请求推进与逐跳数据报双层执行模型。上方事务状态机按需触发报文下沉，下方流水线统一处理每条普通 IP 数据报的逐跳交付。}
  \label{fig:request-and-datagram}
\end{figure}

\subsection{请求推进：已有状态决定哪些步骤根本不会发生}

如果浏览器的 HTTP 缓存已经足够满足当前请求，网络访问可以整体消失。若必须联网，再依次检查：

\begin{itemize}
  \item 主机有没有可用的 IP 地址、前缀、默认网关和 DNS 服务器配置；
  \item URL 中的主机名是否已经对应到可用的目标 IP；
  \item 主场景需要 TCP 时，是否已经有可复用的连接；
  \item 最后才是实际的 HTTP 请求和响应。
\end{itemize}

这里的“依次检查”表示前置条件关系，不表示每次都重新执行这些协议。

\subsection{逐跳交付：当前节点只负责把数据报送到下一节点}

对于一条普通 IP 数据报，当前主机或路由器真正执行的是：
\[
\text{目的 IP}
\xrightarrow{\text{查当前 FIB}}
(\text{出接口},\text{下一跳})
\xrightarrow{\text{准备本跳链路交付}}
\text{下一节点}.
\]

“准备本跳链路交付”不是固定等于 ARP：
\begin{itemize}
  \item Ethernet / 802.11 场景中，若缺少当前下一跳的链路地址，可能需要 ARP 或 ND；
  \item 点对点 PPP 只有一个对端，不需要 Ethernet 式 MAC 查询；
  \item 换到另一种链路时，使用该链路自己的封装和寻址方式。
\end{itemize}

若下一节点仍是路由器，入链路帧在该路由器处结束；IPv4 转发时路由器处理 TTL，再用自己的 FIB 重新决定下一跳。这个过程不断重复，直到目的主机接收数据报。

\section{看综合题时先问四个具体问题}

不需要先给题目贴“这是 DNS 题还是 TCP 题”的标签。先回答下面四问，通常就能知道当前应该调用哪个专题。

\begin{handbooklongtable}{L{2.2cm} Y Y}
\toprule
先问什么 & 要看清什么 & 例子\\
\midrule
现在处理的对象是什么？ & 应用层报文、TCP 报文段、IP 数据报、链路帧还是物理信号 & HTTP 状态码不能拿来解释路由器转发\\
这个机制影响到哪里？ & 当前链路、广播域、一个路由跳、通信两端、一个 AS 或某个应用服务 & ARP 不跨路由器；TCP 连接由两端维护\\
哪台设备或哪个协议模块保存状态？ & 主机、交换机、路由器、路由协议进程、DNS 解析器、TCP 端点等 & MAC 表在交换机；FIB 在路由器；cwnd 在发送端\\
什么事件会让状态改变？ & 报文到达、ACK、超时、缓存过期、路由更新、链路变化等 & 新 LSA 会改变路由知识；ACK 会改变 TCP 发送端状态\\
\bottomrule
\end{handbooklongtable}

控制面与数据面的区别也从这四问自然得到：OSPF/BGP 负责形成和更新路由状态；普通数据报到来时，路由器只使用当前已经安装的 FIB，不会为了这一个包临时重新运行路由协议。

\section{贯穿母场景：先把真实拓扑画对}

主场景固定为 IPv4 + HTTP over TCP（全链路拓扑结构如图~\ref{fig:canonical-topology} 所示）。互相替代的技术不强行塞进同一条真实路径，而放到后面的替换场景中比较。

\begin{figure}[!htbp]
  \centering
  \resizebox{0.92\linewidth}{!}{% ── V04 ISP骨干网与Web服务全流程拓扑 ──
\begin{tikzpicture}[
  scale=1.0,
  >=Stealth,
  font=\small,
  color=themefg,
  text=themefg
]

% ══════════════════════════════════════════════════════════════
% ── 1. 底层连线与区域虚线背景 ──
% ══════════════════════════════════════════════════════════════
\begin{pgfonlayer}{background}
  % Internet / 骨干网区域淡色包络
  \fill[fill=themecurve!6!themebg, draw=themecurve!30, dashed, line width=0.6pt, rounded corners=6pt]
    (-0.5, 3.4) rectangle (7.0, 11.8);
  \node[anchor=north west, font=\footnotesize\bfseries, color=themecurve] at (-0.3, 11.6) {ISP 骨干网络 (Internet WAN)};

  % 公司网络区域淡色包络
  \fill[fill=themeamber!6!themebg, draw=themeamber!30, dashed, line width=0.6pt, rounded corners=6pt]
    (7.6, 3.4) rectangle (19.2, 11.8);
  \node[anchor=north west, font=\footnotesize\bfseries, color=themeamber] at (7.8, 11.6) {某互联网公司网络 (200.1.1.0/24)};

  % Internet 云连线 -> ISP 路由器 E3
  \draw[themegray!80, line width=1.2pt] (1.5, 9.2) -- (3.0, 7.5)
    node[midway, above, sloped, font=\tiny\bfseries, color=themegray] {E3: 全球骨干};

  % ISP 路由器 E0 -> 学校网络
  \draw[themegray!80, line width=1.0pt, ->] (3.6, 8.2) -- (3.6, 10.4)
    node[anchor=south, font=\tiny\bfseries, color=themecurve] {E0: 166.1.0.0/16 (去学校)};

  % ISP 路由器 E2 -> 县里网络
  \draw[themegray!80, line width=1.0pt] (3.6, 6.8) -- (3.6, 4.4)
    node[anchor=north, font=\tiny\bfseries, color=themegray] {E2: 128.14.32.0/24 (去县里)};

  % ISP 路由器 E1 -> 公司路由器 C0
  \draw[themealert!80, line width=1.2pt] (4.3, 7.5) -- (9.4, 7.5)
    node[midway, above, font=\tiny\bfseries, color=themealert, fill=themebg, inner sep=1.5pt, rounded corners=1pt] {E1 $\leftrightarrow$ C0 (200.1.1.254/24)};

  % 公司路由器 C2 -> 专线去学校 (示范跨网直达)
  \draw[themegray!80, line width=1.0pt, dashed] (10.1, 8.2) -- (10.1, 10.0)
    node[anchor=south, font=\tiny\bfseries, color=themegray] {C2: 专线去学校 (200.1.1.2)};

  % 公司路由器 C1 -> 内网交换机
  \draw[themeamber!90, line width=1.2pt] (10.8, 7.5) -- (13.4, 7.5)
    node[midway, above, font=\tiny\bfseries, color=themeamber] {C1: 200.1.1.1 (网关)};

  % 内网交换机 -> Web 服务器, H7, H8
  \draw[themealert!80, line width=1.2pt] (14.2, 7.8) -- (16.7, 8.8);
  \draw[themegray!80, line width=1.0pt] (14.2, 7.5) -- (16.7, 6.8);
  \draw[themegray!80, line width=1.0pt] (14.2, 7.2) -- (16.7, 4.8);
\end{pgfonlayer}

% ══════════════════════════════════════════════════════════════
% ── 2. Internet 云与 ISP 骨干路由器 ──
% ══════════════════════════════════════════════════════════════
% Internet 云图
\begin{scope}[xshift=1.2cm, yshift=9.5cm]
  \begin{pgfonlayer}{background}
    \draw[themecurve!80, fill=themecurve!15!themebg, line width=0.8pt]
      (0, 0) circle (0.45)
      (0.4, 0.2) circle (0.4)
      (0.8, 0) circle (0.45)
      (0.4, -0.2) circle (0.35);
  \end{pgfonlayer}
  \begin{pgfonlayer}{foreground}
    \node[font=\tiny\bfseries, color=themecurve] at (0.4, 0) {Internet};
  \end{pgfonlayer}
\end{scope}

% ISP 骨干路由器
\begin{scope}[xshift=3.6cm, yshift=7.5cm]
  \begin{pgfonlayer}{background}
    \draw[themecurve, fill=themecurve!15!themebg, line width=1.0pt] (0, 0) circle (0.65cm);
  \end{pgfonlayer}
  \begin{pgfonlayer}{foreground}
    \draw[themecurve, line width=1.0pt, ->] (-0.4, 0.25) -- (0.4, -0.25);
    \draw[themecurve, line width=1.0pt, ->] (-0.4, -0.25) -- (0.4, 0.25);
    \draw[themecurve, line width=1.0pt, ->] (0.4, 0.25) -- (-0.4, -0.25);
    \draw[themecurve, line width=1.0pt, ->] (0.4, -0.25) -- (-0.4, 0.25);

    \node[anchor=south, font=\bfseries\footnotesize] at (0, 0.95) {ISP 骨干路由器};

    % 接口标记
    \node[draw=themegray, fill=themebg, inner sep=1.2pt, font=\tiny\bfseries, rounded corners=1pt] at (0, 0.75) {E0};
    \node[draw=themegray, fill=themebg, inner sep=1.2pt, font=\tiny\bfseries, rounded corners=1pt] at (0.75, 0) {E1};
    \node[draw=themegray, fill=themebg, inner sep=1.2pt, font=\tiny\bfseries, rounded corners=1pt] at (0, -0.75) {E2};
    \node[draw=themegray, fill=themebg, inner sep=1.2pt, font=\tiny\bfseries, rounded corners=1pt] at (-0.75, 0) {E3};
  \end{pgfonlayer}
\end{scope}

% ISP 骨干路由转发表 (左下，加宽加高)
\begin{scope}[xshift=-0.2cm, yshift=0.0cm]
  \begin{pgfonlayer}{background}
    \fill[fill=themefg!3!themebg, draw=themecurve!40, line width=0.7pt, rounded corners=6pt]
      (-0.3, -0.3) rectangle (6.2, 3.2);
    \fill[themecurve!18!themebg] (0, 2.0) rectangle (5.8, 2.55);
    \draw[themegray!40, line width=0.5pt] (0, 0) rectangle (5.8, 2.55);
    \draw[themegray!40, line width=0.5pt] (0, 0.5) -- (5.8, 0.5);
    \draw[themegray!40, line width=0.5pt] (0, 1.0) -- (5.8, 1.0);
    \draw[themegray!40, line width=0.5pt] (0, 1.5) -- (5.8, 1.5);
    \draw[themegray!60, line width=0.8pt] (0, 2.0) -- (5.8, 2.0);
    \draw[themegray!40, line width=0.5pt] (2.8, 0) -- (2.8, 2.55);
    \draw[themegray!40, line width=0.5pt] (4.4, 0) -- (4.4, 2.55);
  \end{pgfonlayer}

  \begin{pgfonlayer}{foreground}
    \node[anchor=north, font=\small\bfseries, color=themecurve] at (2.9, 2.95) {ISP 骨干路由转发表};
    \node[font=\scriptsize\bfseries, color=themecurve] at (1.4, 2.27) {目的网络 / 掩码};
    \node[font=\scriptsize\bfseries, color=themecurve] at (3.6, 2.27) {下一跳};
    \node[font=\scriptsize\bfseries, color=themecurve] at (5.1, 2.27) {接口};

    \node[font=\scriptsize, color=themefg] at (1.4, 1.75) {166.1.0.0/16};
    \node[font=\scriptsize, color=themegray] at (3.6, 1.75) {学校路由};
    \node[font=\scriptsize\bfseries, color=themecurve] at (5.1, 1.75) {E0};

    \node[font=\scriptsize, color=themefg] at (1.4, 1.25) {200.1.1.0/24};
    \node[font=\scriptsize, color=themegray] at (3.6, 1.25) {公司路由};
    \node[font=\scriptsize\bfseries, color=themeamber] at (5.1, 1.25) {E1};

    \node[font=\scriptsize, color=themefg] at (1.4, 0.75) {128.14.32.0/24};
    \node[font=\scriptsize, color=themegray] at (3.6, 0.75) {县里路由};
    \node[font=\scriptsize\bfseries, color=themegray] at (5.1, 0.75) {E2};

    \node[font=\scriptsize\bfseries, color=themealert] at (1.4, 0.25) {0.0.0.0/0 (默认)};
    \node[font=\scriptsize, color=themegray] at (3.6, 0.25) {骨干下一跳};
    \node[font=\scriptsize\bfseries, color=themealert] at (5.1, 0.25) {E3};
  \end{pgfonlayer}
\end{scope}

% ══════════════════════════════════════════════════════════════
% ── 3. 公司路由器与公司转发表 ──
% ══════════════════════════════════════════════════════════════
% 公司路由器
\begin{scope}[xshift=10.1cm, yshift=7.5cm]
  \begin{pgfonlayer}{background}
    \draw[themeamber, fill=themeamber!15!themebg, line width=1.0pt] (0, 0) circle (0.65cm);
  \end{pgfonlayer}
  \begin{pgfonlayer}{foreground}
    \draw[themeamber, line width=1.0pt, ->] (-0.4, 0.25) -- (0.4, -0.25);
    \draw[themeamber, line width=1.0pt, ->] (-0.4, -0.25) -- (0.4, 0.25);
    \draw[themeamber, line width=1.0pt, ->] (0.4, 0.25) -- (-0.4, -0.25);
    \draw[themeamber, line width=1.0pt, ->] (0.4, -0.25) -- (-0.4, 0.25);

    \node[anchor=south, font=\bfseries\footnotesize] at (0, 0.95) {公司路由器};

    % 接口标记
    \node[draw=themegray, fill=themebg, inner sep=1.2pt, font=\tiny\bfseries, rounded corners=1pt] at (-0.75, 0) {C0};
    \node[draw=themegray, fill=themebg, inner sep=1.2pt, font=\tiny\bfseries, rounded corners=1pt] at (0.75, 0) {C1};
    \node[draw=themegray, fill=themebg, inner sep=1.2pt, font=\tiny\bfseries, rounded corners=1pt] at (0, 0.75) {C2};
  \end{pgfonlayer}
\end{scope}

% 公司路由转发表 (居中下方，加宽加高)
\begin{scope}[xshift=7.2cm, yshift=0.0cm]
  \begin{pgfonlayer}{background}
    \fill[fill=themefg!3!themebg, draw=themeamber!40, line width=0.7pt, rounded corners=6pt]
      (-0.3, -0.3) rectangle (5.9, 2.9);
    \fill[themeamber!18!themebg] (0, 1.8) rectangle (5.5, 2.35);
    \draw[themegray!40, line width=0.5pt] (0, 0) rectangle (5.5, 2.35);
    \draw[themegray!40, line width=0.5pt] (0, 0.6) -- (5.5, 0.6);
    \draw[themegray!40, line width=0.5pt] (0, 1.2) -- (5.5, 1.2);
    \draw[themegray!60, line width=0.8pt] (0, 1.8) -- (5.5, 1.8);
    \draw[themegray!40, line width=0.5pt] (2.6, 0) -- (2.6, 2.35);
    \draw[themegray!40, line width=0.5pt] (4.2, 0) -- (4.2, 2.35);
  \end{pgfonlayer}

  \begin{pgfonlayer}{foreground}
    \node[anchor=north, font=\small\bfseries, color=themeamber] at (2.75, 2.75) {公司路由转发表};
    \node[font=\scriptsize\bfseries, color=themeamber] at (1.3, 2.08) {目的网络 / 掩码};
    \node[font=\scriptsize\bfseries, color=themeamber] at (3.4, 2.08) {下一跳};
    \node[font=\scriptsize\bfseries, color=themeamber] at (4.85, 2.08) {接口};

    \node[font=\scriptsize, color=themefg] at (1.3, 1.5) {200.1.1.0/24};
    \node[font=\scriptsize, color=themegray] at (3.4, 1.5) {直连内网};
    \node[font=\scriptsize\bfseries, color=themeamber] at (4.85, 1.5) {C1};

    \node[font=\scriptsize, color=themefg] at (1.3, 0.9) {166.1.0.0/16};
    \node[font=\scriptsize, color=themegray] at (3.4, 0.9) {学校专线};
    \node[font=\scriptsize\bfseries, color=themecurve] at (4.85, 0.9) {C2};

    \node[font=\scriptsize\bfseries, color=themealert] at (1.3, 0.3) {0.0.0.0/0 (默认)};
    \node[font=\scriptsize, color=themegray] at (3.4, 0.3) {ISP 路由};
    \node[font=\scriptsize\bfseries, color=themealert] at (4.85, 0.3) {C0};
  \end{pgfonlayer}
\end{scope}

% ══════════════════════════════════════════════════════════════
% ── 4. 公司内网交换机与 Web 服务器群 ──
% ══════════════════════════════════════════════════════════════
% 公司内网交换机
\begin{scope}[xshift=13.8cm, yshift=7.5cm]
  \begin{pgfonlayer}{background}
    \fill[fill=themeamber!12!themebg, draw=themeamber!80, line width=0.8pt, rounded corners=2pt]
      (-0.4, -0.3) rectangle (0.4, 0.3);
  \end{pgfonlayer}
  \begin{pgfonlayer}{foreground}
    \draw[color=themeamber, line width=1.0pt, <->] (-0.25, 0.08) -- (0.25, 0.08);
    \draw[color=themeamber, line width=1.0pt, <->] (-0.25, -0.08) -- (0.25, -0.08);
    \node[anchor=north, font=\tiny] at (0, -0.35) {内网交换机};
  \end{pgfonlayer}
\end{scope}

% DiliDili Web 服务器 (核心服务器机柜)
\begin{scope}[xshift=17.2cm, yshift=8.6cm]
  \begin{pgfonlayer}{background}
    \draw[themealert, fill=themealert!20!themebg, line width=0.8pt, rounded corners=2pt] (-0.5, 0) rectangle (0.5, 0.75);
    \draw[themealert!60, line width=0.5pt] (-0.4, 0.5) -- (0.4, 0.5);
    \draw[themealert!60, line width=0.5pt] (-0.4, 0.25) -- (0.4, 0.25);
    \fill[themealert] (-0.3, 0.62) circle (1.2pt);
    \fill[themegreen] (-0.15, 0.62) circle (1.2pt);
  \end{pgfonlayer}
  \begin{pgfonlayer}{foreground}
    \node[font=\tiny\bfseries, color=themealert] at (0.15, 0.62) {Web};
    \node[anchor=south, font=\tiny, align=center] at (0, 0.8) {
      \textbf{\color{themealert}DiliDili Web 服务器 (HTTP:80)}\\[0.1em]
      \color{themefg}IP: 200.1.1.10/24\\
      \color{themegray}MAC: 33:44:55:66:77:88\\
      \color{themeamber}域名: \texttt{www.dilidili.com}\\
      \color{themegray}默认网关: 200.1.1.1
    };
  \end{pgfonlayer}
\end{scope}

% 主机 H7
\begin{scope}[xshift=17.2cm, yshift=6.6cm]
  \begin{pgfonlayer}{background}
    \draw[themeamber, fill=themeamber!20!themebg, line width=0.8pt, rounded corners=2pt] (-0.45, 0) rectangle (0.45, 0.6);
    \draw[themeamber, line width=0.8pt] (0, 0) -- (0, -0.2);
    \draw[themeamber, line width=0.8pt] (-0.25, -0.2) -- (0.25, -0.2);
  \end{pgfonlayer}
  \begin{pgfonlayer}{foreground}
    \node[font=\bfseries, color=themeamber] at (0, 0.3) {H7};
    \node[anchor=north, font=\tiny, align=center] at (0, -0.25) {
      \textbf{\color{themeamber}主机 H7}\\[0.1em]
      \color{themefg}IP: 200.1.1.20/24\\
      \color{themegray}网关: 200.1.1.1
    };
  \end{pgfonlayer}
\end{scope}

% 主机 H8
\begin{scope}[xshift=17.2cm, yshift=4.6cm]
  \begin{pgfonlayer}{background}
    \draw[themeamber, fill=themeamber!20!themebg, line width=0.8pt, rounded corners=2pt] (-0.45, 0) rectangle (0.45, 0.6);
    \draw[themeamber, line width=0.8pt] (0, 0) -- (0, -0.2);
    \draw[themeamber, line width=0.8pt] (-0.25, -0.2) -- (0.25, -0.2);
  \end{pgfonlayer}
  \begin{pgfonlayer}{foreground}
    \node[font=\bfseries, color=themeamber] at (0, 0.3) {H8};
    \node[anchor=north, font=\tiny, align=center] at (0, -0.25) {
      \textbf{\color{themeamber}主机 H8}\\[0.1em]
      \color{themefg}IP: 200.1.1.30/24\\
      \color{themegray}网关: 200.1.1.1
    };
  \end{pgfonlayer}
\end{scope}

% ══════════════════════════════════════════════════════════════
% ── 5. 底部考点总结卡片（无框通栏，宽裕呼吸感） ──
% ══════════════════════════════════════════════════════════════
\begin{scope}[xshift=-0.5cm, yshift=-3.6cm]
  \begin{pgfonlayer}{background}
    \fill[fill=themefg!4!themebg, draw=themegray!30, line width=0.6pt, rounded corners=4pt]
      (-0.3, -0.4) rectangle (19.4, 2.8);
  \end{pgfonlayer}
  \begin{pgfonlayer}{foreground}
    \node[anchor=north west, text width=19.0cm] at (0.0, 2.45) {
      \small \textbf{\color{themecurve}核心考点精要与机制辨析：}\\[0.4em]
      \scriptsize $\bullet$ \textbf{全流程穿越：一个 Web 请求的一生}：① \textbf{配置接入}（DHCP 获取客户机 IP/掩码/网关/DNS）$\to$ ② \textbf{网关解析}（ARP 广播获取默认网关 MAC）$\to$ ③ \textbf{域名解析}（递归/迭代 DNS 查询获取目标 IP 200.1.1.10）$\to$ ④ \textbf{端点建联}（TCP 三次握手 SYN $\to$ SYN+ACK $\to$ ACK）$\to$ ⑤ \textbf{应用交互}（HTTP GET 请求报文与 200 OK 响应）$\to$ ⑥ \textbf{释放连接}（TCP 四次挥手 FIN/ACK）。\\[0.3em]
      \scriptsize $\bullet$ \textbf{五层协议体系分工与协同全景}：
      \textbf{应用层}（HTTP 报文交互、DNS 名字解析）；
      \textbf{传输层}（TCP 提供端到端可靠字节流、端口多路复用与拥塞控制）；
      \textbf{网络层}（IP 数据报路由转发、CIDR 最长前缀匹配、TTL 逐跳递减）；
      \textbf{数据链路层}（以太网帧封装、MAC 物理寻址、交换机自学习透明转发）；
      \textbf{物理层}（调制与透明比特流传输）。
    };
  \end{pgfonlayer}
\end{scope}

\end{tikzpicture}
}
  \caption{贯穿母场景拓扑结构。全流程跨越局域接入、骨干路由与数据中心三大作用域；实线表示实际数据转发链路，虚线表示 eBGP 控制关系。}
  \label{fig:canonical-topology}
\end{figure}

\subsection{固定示例地址}

\begin{handbooktable}{L{4.0cm} Y}
\toprule
对象 & 示例地址 / 作用\\
\midrule
客户端 & \texttt{192.168.10.23/24}\\
R1 的 LAN 接口 & \texttt{192.168.10.1/24}，作为默认网关\\
R1--R2 点对点 WAN & \texttt{198.51.100.8/30}；R1 为 \texttt{198.51.100.9}，R2 为 \texttt{198.51.100.10}\\
递归 DNS 服务器 & \texttt{192.0.2.53:53}\\
DHCP 服务器 & \texttt{192.0.2.67:67}\\
服务器 LAN 网关 Rn & \texttt{203.0.113.1/24}\\
Web 服务器 & \texttt{203.0.113.80:80}，域名 \texttt{www.example.test}\\
\bottomrule
\end{handbooktable}

这些地址只为了让后面的字段变化可追踪，不表示真实 Internet 地址分配。

\subsection{这张图应该怎样读}

主拓扑只回答“设备在哪里、链路是什么、哪些设备属于同一个作用域”。横向或纵向距离不表示真实距离和时延。

\begin{itemize}
  \item STA、AP、S1、S2 和 R1 的 LAN 接口处在接入网络。S1--S2 有冗余链路，STP 只让其中一条进入普通转发路径。
  \item R1--R2 是 Native PPP 点对点链路，所以这一段没有 Ethernet 源/目的 MAC。
  \item AS64501 内由 OSPF 等域内路由机制维护路由状态；BR1 与 BR2 之间通过 eBGP 交换跨 AS 的前缀和路径属性。
  \item DNS 和 DHCP 服务器挂在 ISP 可达的服务网段。客户端访问它们仍然要经过普通 IP 转发；图中没有画“客户端直连 DNS”的假链路。
\end{itemize}

\subsection{路由表从图里拆出来单独读}

拓扑图不应该同时塞入大块转发表。真正转发时，每台设备只读自己的当前 FIB。例如：

\begin{handbooktable}{L{3.2cm} Y Y}
\toprule
设备 & 典型 FIB 状态 & 当前数据报怎么走\\
\midrule
客户端 & \texttt{192.168.10.0/24} 为直连网段；默认路由经 \texttt{192.168.10.1} & 访问远端 Web 时，IP 目的地址仍是 Web 服务器，但下一跳是 R1\\
R1 & LAN 直连；\texttt{198.51.100.8/30} 直连；其余前缀经上游 R2 & LPM 先确定 PPP 出接口，再做该链路的封装\\
Rn & \texttt{203.0.113.0/24} 直连；其余前缀经上游 & 访问 \texttt{203.0.113.80} 时，下一跳终于就是最终目的主机\\
\bottomrule
\end{handbooktable}

若题目再加入 \texttt{/32}、聚合前缀和默认路由，仍然只做最长前缀匹配：匹配成功的候选中，前缀越长越具体。

\section{同一个 URL 为什么可能产生完全不同的报文}

真正的输入不只有 URL，还包括当前已经保存的状态：

\begin{verbatim}
HostConfig   = 有效 / 缺失 / 过期
DNSCache     = 命中 / 未命中 / 过期
InstalledFIB = 可用 / 缺路由 / 旧版本
Neighbor     = 命中 / 未命中 / 正在解析
SwitchState  = 已学习 / 目的未知 / STP 阻塞
NATState     = 已有映射 / 无映射
TCPState     = 可复用 / 不存在 / 已复位 / 正在关闭
HTTPState    = 可直接使用缓存 / 需重新验证 / 未命中
\end{verbatim}

每个状态都有自己的匹配条件、作用范围和有效时间。“刚访问过一次”不等于所有状态都一定有效。状态剪裁的判定决策流如图~\ref{fig:state-reuse} 所示。

\begin{figure}[!htbp]
  \centering
  \resizebox{0.84\linewidth}{!}{% ── NET-I01 Canonical State Pruning DAG ──
\begin{tikzpicture}[
  >=Stealth,
  font=\small,
  color=themefg,
  text=themefg,
  gate/.style={draw=themecurve,fill=themecurve!9!themebg,rounded corners=4pt,minimum width=3.4cm,minimum height=0.9cm,align=center,font=\scriptsize\bfseries},
  create/.style={draw=themeamber,fill=themeamber!9!themebg,rounded corners=4pt,minimum width=3.0cm,minimum height=0.85cm,align=center,font=\scriptsize\bfseries},
  reuse/.style={draw=themegreen,fill=themegreen!9!themebg,rounded corners=4pt,minimum width=2.8cm,minimum height=0.82cm,align=center,font=\scriptsize\bfseries},
  edge/.style={->,draw=themefg!85,line width=0.85pt},
  bypass/.style={->,draw=themegreen,line width=0.8pt,dashed}
]

\node[gate] (intent) at (7.4,8.0) {用户在浏览器输入 URL};
\node[gate] (httpcache) at (7.4,6.65) {本地是否存在有效强缓存？};
\node[reuse] (localdone) at (11.8,6.65) {直接读取本地缓存渲染};
\draw[edge] (intent) -- (httpcache);
\draw[bypass] (httpcache) -- node[above=2pt,font=\tiny\bfseries,color=themegreen]{命中 (200 from cache)} (localdone);

\node[gate] (config) at (7.4,5.15) {本机是否已配置有效 IP 与网关？};
\node[create] (dhcp) at (2.2,5.15) {触发 DHCP 动态获取配置};
\draw[edge] (httpcache) -- node[right,font=\tiny\bfseries,color=themealert]{未命中 (需联网)} (config);
\draw[edge] (config) -- node[above,font=\tiny\bfseries,color=themealert]{未配置 / 已过期} (dhcp);
\draw[edge] (dhcp.south) |- (7.4,4.45);
\draw[bypass] (config) -- node[right,font=\tiny\bfseries,color=themegreen]{已有配置 (valid)} (7.4,4.45);

\node[gate] (dns) at (7.4,3.75) {是否已知目标服务器 IP 地址？};
\node[create] (resolve) at (2.2,3.75) {发送 DNS 报文进行域名解析};
\draw[edge] (7.4,4.45) -- (dns);
\draw[edge] (dns) -- node[above,font=\tiny\bfseries,color=themealert]{无 DNS 缓存 / 域名未解析} (resolve);
\draw[edge] (resolve.south) |- (7.4,3.05);
\draw[bypass] (dns) -- node[right,font=\tiny\bfseries,color=themegreen]{DNS 缓存命中 / 直填 IP} (7.4,3.05);

\node[gate] (tcp) at (7.4,2.35) {是否存在可复用的 TCP 连接？};
\node[create] (handshake) at (2.2,2.35) {执行 TCP 三次握手建立连接};
\draw[edge] (7.4,3.05) -- (tcp);
\draw[edge] (tcp) -- node[above,font=\tiny\bfseries,color=themealert]{无可用连接} (handshake);
\draw[edge] (handshake.south) |- (7.4,1.65);
\draw[bypass] (tcp) -- node[right,font=\tiny\bfseries,color=themegreen]{可复用 (Keep-Alive)} (7.4,1.65);

\node[gate,draw=themealert,fill=themealert!8!themebg] (http) at (7.4,0.95) {发送 HTTP 请求（或协商缓存 304）并接收响应};
\draw[edge] (7.4,1.65) -- (http);

\begin{pgfonlayer}{background}
  % 左侧创建动作分组框（紧凑包裹 3 个动作节点）
  \fill[fill=themeamber!4!themebg,draw=themeamber!35,dashed,rounded corners=5pt]
    (0.45,1.70) rectangle (3.95,5.75);
  % 右侧强缓存命中终结分组框
  \fill[fill=themegreen!3!themebg,draw=themegreen!30,dashed,rounded corners=5pt]
    (10.15,6.05) rectangle (13.45,7.25);
  % 底部双层执行原则说明卡片
  \fill[fill=themecurve!5!themebg,draw=themecurve!30,rounded corners=4pt]
    (0.45,-0.65) rectangle (13.45,0.25);
\end{pgfonlayer}

\begin{pgfonlayer}{foreground}
  \node[anchor=south,font=\scriptsize\bfseries,color=themeamber] at (2.2,5.80) {仅当状态缺失时才执行创建};
  \node[anchor=center,font=\scriptsize,color=themefg!90] at (6.95,0.0)
    {★ \textbf{冷启动（无缓存/未配置）}走完整创建流程；\textbf{已有状态命中}则直接跳过对应建立步骤。};
  \node[anchor=center,font=\scriptsize,color=themefg!90] at (6.95,-0.38)
    {★ \textbf{DNS 解析、TCP 握手与 HTTP 请求}产生的报文，均走底层逐跳转发通道。};
\end{pgfonlayer}
\end{tikzpicture}
}
  \caption{状态命中与网络事件剪裁 DAG。已有配置、DNS 结果和 TCP 连接可直接复用；各阶段产生的报文均走底层逐跳转发通道。}
  \label{fig:state-reuse}
\end{figure}

这张图还有一个重要用途：它阻止我们把“第二次访问”简单写成“什么都不用做”。例如 DNS 缓存命中只说明不用重新解析名字；它并不能自动保证 TCP 连接仍可复用，也不能保证 ARP 表项还在，更不能保证 HTTP 缓存一定新鲜。

\section{开始普通通信前：没有网络配置时先完成 DHCP}

如果主机已经有有效配置，本节直接跳过。若没有，IPv4 主场景先使用 DHCP 获得：
\[
(\text{主机 IP/前缀},\text{默认网关},\text{DNS 服务器},\text{租约有效期}).
\]

客户端刚启动时还没有普通单播源地址。经典 DORA 过程是 Discover、Offer、Request、ACK；初始报文可以使用源地址 \texttt{0.0.0.0} 和受限广播 \texttt{255.255.255.255}，客户端使用 UDP 68，服务器使用 UDP 67。

若 DHCP 服务器不在本子网，路由器不会直接转发原始二层广播。R1 上的 DHCP Relay（中继）先接收本地 DHCP 消息，再把请求按普通 IP 路由送到远端服务器。这个例子很好地说明：\textbf{应用层协议可以用来建立网络层配置，但协议所在层和它最终修改的状态不必相同。}

IPv6 不照搬这条路径。常见主机可通过 Router Advertisement / SLAAC 获得地址和前缀相关信息，并用 Neighbor Discovery 完成邻居发现；ND 使用 ICMPv6 multicast，不是“IPv6 版 ARP 广播”。

\section{DNS：先判断是否需要解析名字，再把查询报文真正送出去}

浏览器先从 URL 中取出 host（主机部分）。再分三种情况，不要一看到 URL 就默认执行 DNS：

\begin{handbooktable}{L{4.0cm} Y}
\toprule
当前情况 & 这一步实际发生什么\\
\midrule
host 本身就是 IP 地址 & 跳过 DNS，直接把这个 IP 作为后续通信的目标地址\\
DNS 缓存中有仍可使用的记录 & 复用已有结果，不重新发送 DNS 查询\\
host 是主机名，而且没有可用缓存 & 创建一次 DNS 查询，先把主机名解析成目标 IP\\
\bottomrule
\end{handbooktable}

当 DNS 缓存未命中时，客户端已经知道递归 DNS 服务器的 IP，例如 \texttt{192.0.2.53}。但“知道 DNS 服务器在哪”并不等于查询已经到达它。DNS 查询报文本身也必须先被发送出去，客户端仍要完成：
\[
\begin{aligned}
\text{递归 DNS 服务器 IP}
&\xrightarrow{\text{查本机 FIB / LPM}}
(\text{出接口},\text{下一跳})\\
&\xrightarrow{\text{当前链路需要时查邻居表 / ARP}}
\text{本跳帧}\\
&\xrightarrow{\text{逐跳转发}}
\text{递归 DNS 服务器}.
\end{aligned}
\]

因此一个非常重要的顺序是：\textbf{ARP/ND 完全可能发生在目标网站的 DNS 应答之前。}这里 ARP/ND 解析的是“当前 DNS 查询报文的下一跳”，通常是默认网关；它不是在询问 \texttt{www.example.test} 对应 Web server 的 MAC。

递归 DNS 服务器收到查询后，如果自己的缓存也没有答案，它再从自己的网络环境出发，按需要查询根 DNS、顶级域 DNS（TLD）和权威 DNS 服务器。这里发生的是\textbf{递归 DNS 服务器自己的后续查询}；不能把客户端画成亲自依次与这些服务器通信。

DNS 最终交给客户端的是名字解析结果，例如一个或多个 A/AAAA 地址及相应 TTL：
\[
\boxed{\text{domain name}\xrightarrow{\text{DNS}}\text{destination IP candidate(s) + DNS TTL}}
\]
它不输出 next-hop MAC，也不保证这个 Web server 此刻一定可达。CNAME、NS、MX 等记录的完整语义仍由 NET08 负责；本册只关心 DNS 在完整请求中输出什么，以及下一步怎样接到 IP 转发。

\section{核心训练单元：每出现一条 IP 数据报，就跑一次逐跳子程序}

从这里开始，不再为 DNS、TCP 和 HTTP 各背一套“底层流程”。只要题目说“发送这条报文”，先看这条报文形成的 IP 数据报，然后执行同一个逐跳检查：

\begin{methodbox}{逐跳子程序}
\begin{enumerate}
  \item \textbf{写出这条数据报的 destination IP。}DNS query 的 destination 是递归 DNS 服务器；TCP SYN 和后续 HTTP 数据的 destination 才是 Web server。
  \item \textbf{查当前节点的 FIB。}用 LPM 得到出接口和 next-hop IP；若目的位于直连网段，next hop 就是目的主机本身。
  \item \textbf{准备当前这一跳。}Ethernet/WLAN 若缺少下一跳的链路地址，再查 neighbor cache 并按需要执行 ARP/ND；Native PPP 只有一个对端，不做 Ethernet 式 MAC 查询。
  \item \textbf{按当前链路封装并发送。}这一跳使用 802.11、Ethernet、PPP 或题设给出的其他链路表示。
  \item \textbf{如果下一节点是路由器，结束旧的链路帧。}IPv4 路由器处理 TTL 和首部，再从“查自己的 FIB”重新开始；不要把上一跳 MAC 带到下一跳。
  \item \textbf{如果已经到达目的主机，才向 TCP/UDP 交付。}随后由对应的 TCP/UDP 模块继续处理端口、序号、ACK 或应用数据。
\end{enumerate}
\end{methodbox}

这个子程序不是说每条链路做完全相同的事，而是说\textbf{每一跳都要重新回答同一组问题：最终目的 IP 是谁、当前 FIB 让我从哪里出去、这一跳先交给谁、当前链路怎样把数据送过去。}

下面选一条从客户端发往 Web 服务器的 IPv4 数据报，按设备顺序把这个子程序完整跑一遍。DNS 查询、TCP SYN、TCP ACK 和承载 HTTP 数据的报文段都会重复调用它；只是它们的目的 IP、方向和上层语义可能不同。

\subsection{客户端：先判断目标是不是本地，再确定下一跳}

目标地址为 \texttt{203.0.113.80}。客户端根据自己的前缀判断它不在 \texttt{192.168.10.0/24} 内，于是查 FIB 命中默认路由：
\[
\text{destination IP}=\texttt{203.0.113.80},
\qquad
\text{next-hop IP}=\texttt{192.168.10.1}.
\]

这两个“目的”必须同时保留：IP 层描述的最终目的仍是 Web 服务器；当前链路只负责先把帧交给 R1。

如果 ARP cache 中没有 R1 的 MAC，客户端就在 VLAN 10 的当前广播域内发送 ARP Request。ARP Request 不会穿过 R1 去询问远端 Web server。

\subsection{STA 到 AP：802.11 只解决这一段无线链路}

经典 infrastructure BSS 中，STA 向 distribution system 发送数据时 \texttt{To DS=1, From DS=0}：

\begin{itemize}
  \item Address 1：RA/BSSID，即当前无线接收端 AP；
  \item Address 2：TA/SA，即发送 STA；
  \item Address 3：distribution system 一侧的最终二层目的，本例为 R1 的 LAN MAC；
  \item Address 4：这一方向不使用。只有 \texttt{To DS=1, From DS=1} 的无线桥接/WDS 等场景才需要第四地址。
\end{itemize}

CSMA/CA、DIFS/SIFS、随机退避、NAV、RTS/CTS 和 MAC ACK 都只服务当前 WLAN hop。收到 802.11 ACK 只能说明这一跳得到链路层反馈，不能推出远端 HTTP 已成功。

\subsection{AP 到交换网络：换成 Ethernet 表示，但还没有经过 IP router}

AP 在主场景中是 bridge。它把 802.11 数据帧中的上层 packet 转为 Ethernet 帧继续送入 VLAN 10：

\begin{itemize}
  \item Ethernet source MAC 仍对应 STA；
  \item Ethernet destination MAC 为当前三层下一跳 R1；
  \item 无标签 Ethernet II 中 IPv4 EtherType 为 \texttt{0x0800}；
  \item 若在 trunk 上使用 802.1Q，则会插入 VLAN tag，常见 TPID 为 \texttt{0x8100}；
  \item FCS 属于当前帧的检错信息，不是端到端校验。
\end{itemize}

AP bridge 和二层交换机都不是 IP router，因此不会因为经过它们就让 IPv4 TTL 减 1。

\subsection{Switch：先学习 source，再根据 destination 决定怎么转发}

交换机收到 frame 后，先用 source MAC 更新 MAC table，再查 destination MAC：

\begin{enumerate}
  \item destination 未知：在同 VLAN、且 STP 允许转发的其他端口上泛洪；
  \item destination 已知且在另一个端口：定向转发；
  \item destination 已知且就在入端口：过滤，不再发回该端口；
  \item broadcast/multicast：按相应二层规则处理。
\end{enumerate}

STP 解决的是冗余二层链路产生环路的问题。它先形成一套 active topology；普通 frame 到来时只使用当前端口状态，不是每个 frame 都重新运行一次 STP。

\subsection{R1：这里才真正进入下一次 IP hop}

R1 接收发给自己 LAN MAC 的 Ethernet frame，验证并结束当前链路封装，然后处理 IPv4 packet：

\begin{enumerate}
  \item 检查并更新 IPv4 首部相关状态；
  \item TTL 减 1；若结果为 0，丢弃原 packet 并产生 ICMP Time Exceeded；
  \item 对 destination IP 做 LPM，决定新的出接口和下一跳；
  \item 按新的出接口建立下一段链路封装。
\end{enumerate}

本例 R1 还执行 NAPT。私网端点
\[
\texttt{192.168.10.23}:p
\]
可映射为公网端点
\[
\texttt{198.51.100.9}:p'.
\]
因此外部看到的 source IP 和 source port 都可能变化。IPv4 首部校验和必须相应更新；TCP/UDP 校验和还覆盖 pseudo-header 中的 IP 信息，所以受影响的传输层校验和也要调整。回程数据必须命中反向 NAT mapping 才能回到内部 client。

NAPT 只负责地址/端口映射，并不负责 TCP 的可靠传输。Seq/Ack、重传和窗口状态仍由两端的 TCP 模块维护。

\subsection{R1 到 R2：Native PPP 与 PPPoE 不能混在一起}

R1--R2 是 Native PPP 点对点链路。这里只有一个 peer，所以不需要 Ethernet 式目的/源 MAC 学习。

若题目继续追问 PPP 帧，经典 HDLC-like PPP 帧依次包含 Flag、Address、Control、Protocol、Information、FCS 和尾部 Flag。Address/Control 不是 Ethernet 的源/目的 MAC；异步链路常用字节转义，比特同步链路使用零比特填充。LCP 用于建立和配置 PPP 链路，NCP 再配置对应网络层协议。经典 PPP 可以检错并丢弃坏帧，但不因此拥有 TCP 式序号、ACK 和可靠重传。

PPPoE 则不同：它把 PPP 放在 Ethernet 上承载，所以外层仍然存在 Ethernet MAC。看到 PPPoE 不能写成“PPP 所以这段没有 MAC”。

\subsection{中间路由器：重复同一个动作，不为当前 packet 临时跑一次 OSPF/BGP}

R2、Core、BR1、BR2、Rn 都只对当前 packet 使用自己已经安装的转发表：

\[
(\text{destination IP},\text{current FIB})
\mapsto
(\text{出接口},\text{下一跳}).
\]

每经过一个真正的 IP router，TTL 再减 1，并按新出接口重新准备本跳链路交付。OSPF/BGP 的职责是更慢时间尺度上的路由信息交换、路径选择和 FIB 更新，不是 packet 到来以后才开始计算。

\subsection{最后一跳：这时 next hop 才等于最终 Web server}

Rn 的 FIB 命中直连前缀 \texttt{203.0.113.0/24}，因此：
\[
\text{next-hop IP}=\text{destination IP}=\texttt{203.0.113.80}.
\]
如果 ARP cache 未命中，Rn 对 Web server 做 ARP，得到 Web MAC，再创建最后一跳 Ethernet frame。服务器网卡接收后，IPv4 把数据交给相应的 TCP/UDP 模块。

\section{把字段放在正确的作用范围里：哪些保持，哪些会变}

前面的逐设备解释很长，可以压成一张字段与状态演进全景图（如图~\ref{fig:field-lifetime} 所示）。读图时不要用一句“包一路不变”或“每跳全都变”概括，因为不同层的字段寿命并不一样。

\begin{figure}[!htbp]
  \centering
  \resizebox{0.92\linewidth}{!}{% ── NET-I01 Canonical Cross-Media Frame & State Evolution Diagram ──
\begin{tikzpicture}[
  scale=1.0,
  >=Stealth,
  font=\small,
  color=themefg,
  text=themefg,
  stagehdr/.style={draw=themefg!70,fill=themefg!8!themebg,rounded corners=4pt,minimum width=3.4cm,minimum height=1.05cm,align=center,font=\scriptsize\bfseries,line width=1.0pt},
  hdrfield/.style={draw=themegreen!80,fill=themegreen!12!themebg,rounded corners=2pt,font=\fontsize{5.2pt}{6.2pt}\selectfont,minimum height=1.2cm,inner sep=2pt,align=left},
  payloadfield/.style={draw=themecurve,fill=themecurve!18!themebg,rounded corners=2pt,font=\fontsize{5.5pt}{6.5pt}\selectfont\bfseries,minimum height=1.2cm,inner sep=2pt,align=center},
  fcsfield/.style={draw=themealert,fill=themealert!18!themebg,rounded corners=2pt,font=\fontsize{5.2pt}{6.2pt}\selectfont\bfseries,minimum height=1.2cm,inner sep=2pt,align=center}
]

\def\colA{1.8}
\def\colB{5.6}
\def\colC{9.4}
\def\colD{13.2}
\def\colE{17.0}

\begin{pgfonlayer}{background}
  \fill[fill=themeamber!6!themebg,draw=themeamber!45,dashed,line width=0.9pt,rounded corners=8pt]
    (-0.3, 10.4) rectangle (19.0, 13.8);
  \fill[fill=themecurve!6!themebg,draw=themecurve!45,dashed,line width=0.9pt,rounded corners=8pt]
    (-0.3, 6.4) rectangle (19.0, 10.0);
  \fill[fill=themegreen!6!themebg,draw=themegreen!45,dashed,line width=0.9pt,rounded corners=8pt]
    (-0.3, 1.4) rectangle (19.0, 6.0);
  \fill[fill=themecurve!8!themebg,draw=themecurve!50,line width=0.9pt,rounded corners=6pt]
    (-0.3, -1.1) rectangle (19.0, 1.0);
  \draw[themefg!15,dashed,line width=0.6pt] (3.7, 1.4) -- (3.7, 13.8);
  \draw[themefg!15,dashed,line width=0.6pt] (7.5, 1.4) -- (7.5, 13.8);
  \draw[themefg!15,dashed,line width=0.6pt] (11.3, 1.4) -- (11.3, 13.8);
  \draw[themefg!15,dashed,line width=0.6pt] (15.1, 1.4) -- (15.1, 13.8);
\end{pgfonlayer}

\begin{pgfonlayer}{foreground}
  \node[anchor=west,font=\small\bfseries,color=themeamber] at (-0.1, 13.4) {【TCP】序号、确认号和可靠传输状态由通信两端维护};
  \node[anchor=west,font=\small\bfseries,color=themecurve] at (-0.1, 9.6) {【IPv4】路由器逐跳转发；TTL 在每个 IP 路由器处递减};
  \node[anchor=west,font=\small\bfseries,color=themegreen] at (-0.1, 5.6) {【链路层】802.11、Ethernet、PPP 使用各自的本跳帧表示};
\end{pgfonlayer}

\node[stagehdr,draw=themegreen] at (\colA, 14.8) {① 主机 $\to$ AP 无线跳\\\tiny 802.11 CSMA/CA 信道};
\node[stagehdr,draw=themegreen] at (\colB, 14.8) {② AP $\to$ 交换机 $\to$ 网关 R1\\\tiny 802.3 双绞线以太网};
\node[stagehdr,draw=themecurve] at (\colC, 14.8) {③ R1 $\to$ R2 广域网跳\\\tiny Native PPP 点对点光纤};
\node[stagehdr,draw=themecurve] at (\colD, 14.8) {④ 运营商 / 跨 AS 路由器\\\tiny 按已安装 FIB 逐跳转发};
\node[stagehdr,draw=themeamber] at (\colE, 14.8) {⑤ 网关 Rn $\to$ Web 服务器\\\tiny 服务器局域网以太网};

\draw[->,draw=themefg!80,line width=1.3pt] (\colA+1.8, 14.8) -- (\colB-1.8, 14.8);
\draw[->,draw=themefg!80,line width=1.3pt] (\colB+1.8, 14.8) -- (\colC-1.8, 14.8);
\draw[->,draw=themefg!80,line width=1.3pt] (\colC+1.8, 14.8) -- (\colD-1.8, 14.8);
\draw[->,draw=themefg!80,line width=1.3pt] (\colD+1.8, 14.8) -- (\colE-1.8, 14.8);

\node[align=center,font=\scriptsize] at (\colA, 12.0) {
  \textbf{\color{themeamber}TCP 端点:}\\[1pt]
  \texttt{src=192.168.10.23:$p$}\\
  \texttt{dst=203.0.113.80:80}\\[2pt]
  \textbf{\color{themefg}Seq / Ack:} 由两端 TCP 维护
};

\node[align=center,font=\scriptsize] at (\colB, 12.0) {
  \textbf{\color{themeamber}TCP 端点仍是两端主机:}\\[1pt]
  \texttt{src=192.168.10.23:$p$}\\
  \texttt{dst=203.0.113.80:80}\\[2pt]
  \textbf{\color{themegray}AP / switch 不维护 TCP 状态}
};

\node[align=center,font=\scriptsize] at (\colC, 12.0) {
  \textbf{\color{themealert}R1 执行 NAPT:}\\
  \texttt{src=198.51.100.9:$p'$}\\
  \texttt{dst=203.0.113.80:80}\\[2pt]
  \textbf{\color{themefg}Seq / Ack 仍由两端 TCP 维护}
};

\node[align=center,font=\scriptsize] at (\colD, 12.0) {
  \textbf{\color{themeamber}公网侧看到的端点:}\\[1pt]
  \texttt{src=198.51.100.9:$p'$}\\
  \texttt{dst=203.0.113.80:80}\\[2pt]
  \textbf{\color{themegray}中间路由器不维护}\\[1pt]
  \textbf{\color{themegray}TCP 连接状态}
};

\node[align=center,font=\scriptsize] at (\colE, 12.0) {
  \textbf{\color{themeamber}到达服务器 TCP:}\\[1pt]
  \texttt{src=198.51.100.9:$p'$}\\
  \texttt{dst=203.0.113.80:80}\\[2pt]
  \textbf{\color{themefg}TCP 处理 Seq / Ack}\\[1pt]
  \textbf{\color{themefg}再把重组后的字节交给 Web 进程}
};

\node[align=center,font=\scriptsize] at (\colA, 8.1) {
  \textbf{\color{themecurve}源 IP:} \texttt{192.168.10.23} (私网)\\[2pt]
  \textbf{\color{themecurve}目的 IP:} \texttt{203.0.113.80} (公网)\\[3pt]
  \textbf{\color{themeamber}TTL 初始值:} $t$ (如 64 或 128)
};

\node[align=center,font=\scriptsize] at (\colB, 8.1) {
  \textbf{\color{themecurve}源 IP 与目的 IP 均不变}\\[2pt]
  \textbf{\color{themegreen}TTL 不减 1!} (二层交换不耗 TTL)\\[3pt]
  \textbf{\color{themegray}首部校验和保持原值}
};

\node[align=center,font=\scriptsize] at (\colC, 8.1) {
  \textbf{\color{themealert}网关 R1 转发动作:}\\
  1. 源 IP 转换为公网 \texttt{198.51.100.9}\\[2pt]
  2. \textbf{\color{themealert}TTL 减 1} $\to t-1$\\[2pt]
  3. \textbf{\color{themealert}重新计算 IP 首部校验和}
};

\node[align=center,font=\scriptsize] at (\colD, 8.1) {
  \textbf{\color{themecurve}目的 IP 203.0.113.80 保持}\\[2pt]
  \textbf{\color{themealert}每过一个 IP 路由器:}\\
  $\text{TTL} \leftarrow \text{TTL} - 1$\\[2pt]
  重新计算 IP 首部校验和
};

\node[align=center,font=\scriptsize] at (\colE, 8.1) {
  \textbf{\color{themecurve}成功到达目标:} \texttt{203.0.113.80}\\[2pt]
  \textbf{\color{themecurve}最终 TTL:} $t - \text{总路由跳数}$\\[3pt]
  Server 剥离 IP 首部交付 TCP
};

\node[align=center] at (\colA, 3.6) {
  \textbf{\color{themegreen}\scriptsize 802.11 无线数据帧}\\[3pt]
  \begin{tikzpicture}
    \node[hdrfield,minimum width=2.05cm] at (0,0) {
      \textbf{802.11 帧首部 (三地址)}\\
      \texttt{RA = AP\_MAC} (接收端)\\
      \texttt{TA = STA\_MAC} (发送端)\\
      \texttt{Addr3 = R1\_MAC} (网关)
    };
    \node[payloadfield,minimum width=0.75cm] at (1.46,0) {IP\\报文};
    \node[fcsfield,minimum width=0.55cm] at (2.16,0) {FCS\\(CRC)};
  \end{tikzpicture}\\[3pt]
  \tiny \color{themegray} 无线网卡硬件计算并追加 FCS
};

\node[align=center] at (\colB, 3.6) {
  \textbf{\color{themegreen}\scriptsize Ethernet II（DIX）帧}\\[3pt]
  \begin{tikzpicture}
    \node[hdrfield,minimum width=2.05cm] at (0,0) {
      \textbf{标准以太网帧首部}\\
      \texttt{目的: R1\_MAC} (网关)\\
      \texttt{源: STA\_MAC} (主机)\\
      \texttt{类型: 0x0800} (IPv4)
    };
    \node[payloadfield,minimum width=0.75cm] at (1.46,0) {IP\\报文};
    \node[fcsfield,minimum width=0.55cm] at (2.16,0) {FCS\\(重算)};
  \end{tikzpicture}\\[3pt]
  \tiny \color{themegray} AP 桥接：剥离 802.11 重构以太网帧
};

\node[align=center] at (\colC, 3.6) {
  \textbf{\color{themecurve}\scriptsize Native PPP 点对点帧}\\[3pt]
  \begin{tikzpicture}
    \node[hdrfield,minimum width=2.05cm,draw=themecurve!80,fill=themecurve!12!themebg] at (0,0) {
      \textbf{PPP 首部 (点对点链路)}\\
      \texttt{Flag = 0x7E, Addr = 0xFF}\\
      \texttt{Control = 0x03}\\
      \texttt{协议 = 0x0021} (IPv4)
    };
    \node[payloadfield,minimum width=0.75cm] at (1.46,0) {IP\\报文};
    \node[fcsfield,minimum width=0.55cm] at (2.16,0) {FCS\\(CRC)};
  \end{tikzpicture}\\[3pt]
  \tiny \bfseries \color{themealert} 点对点链路：无以太网 MAC 地址!
};

\node[align=center] at (\colD, 3.6) {
  \textbf{\color{themecurve}\scriptsize 路由器按出接口重新封装}\\[3pt]
  \begin{tikzpicture}
    \node[hdrfield,minimum width=2.05cm,draw=themecurve!80,fill=themecurve!12!themebg] at (0,0) {
      \textbf{出接口链路首部}\\
      依出接口链路类型确定\\
      PPP / Ethernet / 其他链路\\
      \texttt{建立本跳所需链路信息}
    };
    \node[payloadfield,minimum width=0.75cm] at (1.46,0) {IP\\报文};
    \node[fcsfield,minimum width=0.55cm] at (2.16,0) {FCS\\(每跳)};
  \end{tikzpicture}\\[3pt]
  \tiny \color{themegray} 剥离旧 L2 首部 $\to$ 查表 $\to$ 封装新 L2
};

\node[align=center] at (\colE, 3.6) {
  \textbf{\color{themeamber}\scriptsize 服务器局域网以太网帧}\\[3pt]
  \begin{tikzpicture}
    \node[hdrfield,minimum width=2.05cm,draw=themeamber!80,fill=themeamber!12!themebg] at (0,0) {
      \textbf{服务器以太网帧首部}\\
      \texttt{目的: Web\_MAC} (服务器)\\
      \texttt{源: Rn\_MAC} (网关)\\
      \texttt{类型: 0x0800} (IPv4)
    };
    \node[payloadfield,minimum width=0.75cm] at (1.46,0) {IP\\报文};
    \node[fcsfield,minimum width=0.55cm] at (2.16,0) {FCS\\(校验)};
  \end{tikzpicture}\\[3pt]
  \tiny \bfseries \color{themeamber} 网关 Rn 通过 ARP 解析获得 Web\_MAC
};

\node[anchor=west,font=\small\bfseries,color=themecurve] at (-0.1, 0.70) {★ 读这张图时只分清三件事：};
\node[anchor=west,font=\scriptsize\bfseries,color=themefg] at (-0.1, 0.25) {
  1. \textbf{TCP：}Seq / Ack 和可靠传输状态由两端维护；NAPT 可以改写公网侧看到的源 IP / Port；
};
\node[anchor=west,font=\scriptsize\bfseries,color=themefg] at (-0.1, -0.20) {
  2. \textbf{IPv4：}本例的目的 IP 始终指向 Web server；每经过一个 IP 路由器，TTL 减 1 并更新 IPv4 首部校验和；
};
\node[anchor=west,font=\scriptsize\bfseries,color=themefg] at (-0.1, -0.65) {
  3. \textbf{链路层：}AP 可在 802.11 / Ethernet 间转换表示；经过 IP 路由器后，下一条链路按自己的格式重新封装。
};

\end{tikzpicture}
}
  \caption{跨介质全链路帧转换与端到端状态演进。TCP 状态由两端维护，IP 首部由路由器递减 TTL，链路层帧在每一跳按当前介质重新封装。}
  \label{fig:field-lifetime}
\end{figure}

\begin{handbooklongtable}{L{3.0cm} L{3.0cm} Y}
\toprule
对象 / 字段 & 主要作用范围 & 主场景中怎样变化\\
\midrule
HTTP 资源 / 资源表示 & 应用层两端 & 中间交换机和路由器不解释 HTTP 资源语义\\
TCP Seq/Ack & TCP 两端 & 普通路由器和基础 NAPT 不维护 TCP 可靠传输状态\\
TCP/UDP 端口 & 主机内的传输端点 & NAPT 可以改写公网侧看到的端口\\
IPv4 目的地址 & 跨多个路由跳 & 主场景中始终指向 Web server；若题设加入 DNAT/隧道再另行判断\\
IPv4 源地址 & 跨多个路由跳 & 在 R1 的 NAPT 处从私网地址改成公网地址\\
TTL & 每个 IP 路由跳 & 只在 IP 路由器处递减；AP/交换机不消耗 TTL\\
MAC / 802.11 地址字段 & 当前二层交付范围 & AP 可在 802.11 与 Ethernet 之间转换表示；经过路由器后按下一链路重新确定\\
FCS & 当前链路帧 & 每个新帧独立计算和检查，不是端到端状态\\
物理信号 & 当前传输介质 & WLAN、电缆、光纤可以使用完全不同的信号表示\\
\bottomrule
\end{handbooklongtable}

一个重要边界是：\textbf{交换机内部的转发和路由器的逐跳转发不是同一件事。}在同一 Ethernet 二层交付范围内，交换机通常保持源/目的 MAC 的语义；经过 IP router 后，入链路的二层目的已经完成任务，下一条链路再按新的 next hop 建立新的封装。

\section{TCP：把多条 IP 数据报组织成端到端字节流}

主场景中的 HTTP 使用 TCP。TCP 连接真正由两端 TCP 栈维护，中间普通路由器不保存它的 Seq/Ack/window 状态。

\subsection{端口和连接标识}

端口解决的是“这台主机上的哪一个进程/传输端点应该接收数据”。IANA 常把端口范围分为 System Ports \texttt{0--1023}、User Ports \texttt{1024--49151} 和 Dynamic/Private Ports \texttt{49152--65535}；端口 0 保留。

网络题中常用五元组
\[
(\text{protocol},\text{local IP},\text{local port},\text{remote IP},\text{remote port})
\]
描述一条传输流。对已经建立的 TCP 连接，协议已经固定后，四个端点字段足以区分连接。不要把“5-tuple 唯一定位所有 socket”当成无条件定理：监听 socket 和 UDP socket 的分用条件并不完全相同。

\subsection{三次握手和连接释放}

经典三次握手可以写成：
\[
\mathrm{SYN}(x)
\to
\mathrm{SYN+ACK}(y,x+1)
\to
\mathrm{ACK}(x+1,y+1).
\]
每一个箭头都代表一条需要经过前述逐跳网络的 TCP segment / IP datagram。

连接关闭时，FIN/ACK、半关闭、TIME\_WAIT 与异常 RST 都属于 NET06 的 TCP 生命周期。本册主线停在收到目标 HTTP 内容，因此不在这里展开完整挥手状态机；题目若继续追问连接怎样结束，再转回 NET06。

\subsection{可靠性、流量控制和拥塞控制不能混成一个窗口}

TCP sender 同时受到两类不同限制：

\begin{itemize}
  \item receiver 用 \texttt{rwnd} 告诉 sender 自己还能接收多少；
  \item sender 用 \texttt{cwnd} 控制自己愿意向共享网络放入多少未确认数据。
\end{itemize}

若 \texttt{FlightSize} 表示已经发出但尚未确认的数据量，则当前还能新增发送的上界为
\[
W_{\mathrm{usable}}
=
\max\!\left(0,\min(rwnd,cwnd)-FlightSize\right).
\]

\texttt{rwnd=0} 说明接收端暂时没有空间，sender 需要使用零窗口探测等机制等待窗口重新打开；loss、dupACK、RTO 或 ECN 等反馈则可能改变拥塞控制状态。慢开始、拥塞避免、快重传和快恢复由 NET07 完整拥有。

GBN/SR 的窗口、序号空间和回退规则由 NET03 解释；TCP 是更复杂的端到端实例。BDP 则回答“链路中要有多少在途数据才能把带宽充分利用”，通过 NET-B05 与 window 联系起来。

\subsection{UDP 在这条主线里的位置}

DNS 常用 UDP；UDP header 固定包含 source port、destination port、length 和 checksum。其 checksum 计算还会使用包含源/目的 IP、协议号和长度等信息的 pseudo-header，所以 NAT 改写地址或端口时不能只改表面字段而不处理校验和。

UDP 不建立 TCP 那样的连接状态，也不提供 TCP 式可靠字节流。若题目比较 TCP 与 UDP，应回到 NET06，而不是从“是否有端口”判断是否面向连接。

\section{HTTP：字节到了以后，应用层才重新解释“资源”}

TCP 只提供字节流，不知道这些字节代表 HTML、图片还是状态码。HTTP 才解释 method、target、header fields、status 和 resource semantics。

\subsection{资源和资源表示要分开}

\textbf{Resource（资源）}是 URI 所标识的逻辑对象；\textbf{representation（资源表示）}是服务器在某次响应中用来表示这个资源的具体内容和元数据。同一个资源可以有不同表示，因此不能把“资源”直接等同于某一份固定字节文件。

常见 method（方法）如 GET、HEAD、POST、PUT、DELETE、OPTIONS；状态码按 1xx--5xx 分组。本册不重新列完整状态码表，只强调它们属于应用层语义，不能拿 HTTP 状态码替代 TCP/IP 层的成功证据。

\subsection{一个网页往往不是一个对象}

主 HTML 返回后，浏览器才可能从中发现 CSS、JavaScript、图片等 URL。因此页面加载更像一个依赖图：
\[
\text{main HTML}
\to
\text{parse}
\to
\{\text{CSS},\text{JS},\text{images},\ldots\}.
\]

这就是为什么性能题不能只写“HTTP 需要 1 RTT”。必须先看哪些对象必须串行发现，哪些连接能复用，哪些请求可以并行或多路复用（multiplex）。

\subsection{HTTP cache 会改变整个网络事件图}

如果本地缓存仍然新鲜，并且足以直接满足当前请求，后续网络交互可以整体消失；若缓存需要重新验证，浏览器可以发送条件请求。服务器返回 304 时，表示客户端可以继续使用本地已有的资源表示，而不是“这个请求叫 304”。

\subsection{HTTP 版本改变的是连接与并发组织方式}

\begin{handbooktable}{L{2.3cm} Y Y}
\toprule
版本 & 典型考试观察点 & 不能误读为\\
\midrule
HTTP/1.0 & 默认非持久连接（non-persistent）时，多个对象常反复建立 TCP 连接 & 每个页面永远只有一条连接\\
HTTP/1.1 & 默认持久连接（persistent）；是否允许 pipeline 取决于题设 & 持久连接不等于所有对象自动并行\\
HTTP/2 & 多个 stream 可复用一条 TCP connection & TCP 丢包仍可能使多个 stream 一起等待底层有序字节流\\
HTTP/3 & 基于 QUIC 的独立 stream 减少共享 TCP 字节流造成的跨 stream 阻塞 & 单个 stream 内仍要求可靠、有序交付\\
\bottomrule
\end{handbooktable}

\section{路由协议在后台工作，当前数据报只看当前 FIB}

路由与转发必须分开：
\[
\boxed{\text{路由协议负责形成/更新转发表}}
\qquad
\boxed{\text{当前数据报只使用已经生效的转发表}}.
\]

\begin{handbooklongtable}{L{2.4cm} Y Y}
\toprule
机制 & 它在后台做什么 & 当前数据报真正使用什么\\
\midrule
STP & 在冗余二层拓扑中形成无环的 active topology & switch 当前端口状态\\
Switch learning & 从 source MAC 学习 MAC--port mapping & destination MAC lookup\\
RIP & 距离向量交换，经典 metric 为 hop count，存在 count-to-infinity 等边界 & 安装后的 FIB\\
OSPF & 邻接、LSA、LSDB、SPF 计算域内路由 & 安装后的 FIB\\
BGP & AS 间交换 prefix 和 path attributes，并按 policy 选择路径 & border router 当前已安装的下一跳/出接口\\
\bottomrule
\end{handbooklongtable}

特别注意时间顺序：收到 routing update，不等于新 route 已经 selected；selected route 不等于 FIB 已安装；FIB 已安装也不能反过来改变更早到达的数据报走过的路径。

\section{三类异常不要混在一起：生命周期、尺寸和反馈}

\subsection{TTL / ICMP：数据报还能不能继续往前走}

每经过一个 IP router，IPv4 TTL 减 1。减到 0 的 router 丢弃原 packet，并产生 ICMP Time Exceeded。ICMP Echo Request/Reply 则可以用于主动探测 IP 层可达性。

ICMP 是网络层反馈，不等于 TCP ACK，也不等于“网页一定可以正常访问”。

\subsection{MTU：下一条链路能不能装下这个 IP packet}

链路 MTU 会向上限制 IP packet 的大小：
\[
\text{Link MTU}
\to
\text{IP packet size}
\to
\text{fragmentation / PMTU}
\to
\text{TCP MSS / segment size}.
\]

主场景故意让 PPP WAN 的 MTU 为 800 B。IPv4 packet 超过出接口 MTU 时，是否允许分片要看 DF 和题设；fragment offset、MF 和目的端重组由 NET04 负责。若 TCP 已根据 MSS/PMTU 把 segment 控制在更合适的大小，可以避免不必要的 IP fragmentation。

\subsection{链路 ACK、TCP ACK、rwnd、cwnd、ICMP、HTTP status 各自说明不同事情}

\begin{handbooktable}{L{3.0cm} Y Y}
\toprule
看到什么反馈 & 它主要说明什么 & 不能直接推出什么\\
\midrule
802.11 MAC ACK & 当前无线单跳得到链路层确认 & Web server 已收到 HTTP\\
TCP ACK / dupACK / RTO & TCP byte delivery 的端到端证据与恢复触发 & 接收缓冲区或路由状态本身\\
\texttt{rwnd} & receiver 还能接收多少 & 网络路径是否拥塞\\
\texttt{cwnd} 的变化 & sender 对共享路径拥塞程度的控制结果 & HTTP cache 是否命中\\
ICMP error & IP 层路径/packet 处理出现特定问题 & TCP 已正常关闭\\
HTTP status & application 对 request 的语义响应 & 下层从未发生丢包/重传\\
\bottomrule
\end{handbooktable}

\section{性能题先画依赖关系，再算 RTT 和传输时间}

网络时延的基本账仍来自 NET01：
\[
d_{total}=d_{proc}+d_{queue}+d_{trans}+d_{prop},
\qquad
d_{trans}=\frac{L}{R}.
\]

传播时延由距离和传播速度决定，发送时延由数据量和链路速率决定，二者不能互换。BDP
\[
\mathrm{BDP}=\text{bandwidth}\times\mathrm{RTT}
\]
表示路径上为了充分利用带宽，大致需要容纳多少在途数据。

对于 Web 访问，先画真正的依赖关系：
\[
\text{network configuration if needed}
\to
\text{DNS if needed}
\to
\text{TCP setup if needed}
\to
\text{main HTTP object}
\to
\text{subresource dependencies}.
\]

然后再删掉已经命中的状态，标出能并行的部分，最后才相加 RTT 和 transmission time。不能机械套“DNS 1 RTT + TCP 1 RTT + HTTP 1 RTT”。

\section{如果把主场景的一段网络换掉，应该切换到什么机制}

主场景不需要把所有协议同时运行，但考试可能把其中某一段换掉。此时只替换对应机制，不要让无关层一起变化。

\begin{handbooklongtable}{L{3.2cm} Y Y}
\toprule
题设怎样改 & 应切换到什么机制 & 仍然保持什么\\
\midrule
switched full-duplex Ethernet 改为共享半双工 Ethernet & CSMA/CD、$2\tau$、64 B 最小帧与二进制指数退避的经典模型 & IP destination、TCP 端点语义不变\\
WLAN 出现 hidden station & RTS/CTS、NAV、SIFS/DIFS 更重要 & 仍只解决当前 WLAN hop 的介质访问\\
随机竞争改成 polling / token & 发送权由轮询或令牌控制，会引入等待与控制开销 & 不取代 IP routing\\
共享资源改为 FDM/TDM/CDM & 静态划分频率、时间或正交码；CDM 用码向量相关/内积区分用户 & 不取代链路以上协议\\
Native PPP 改为 PPPoE & PPP 外面再套 Ethernet，MAC 地址重新出现 & PPP 会话/控制语义仍存在\\
IPv4 改为 IPv6 & SLAAC/RA、ND、Hop Limit 等替换 IPv4 对应机制；router 不做 IPv4 式中途分片 & destination 与 next hop 仍要分清\\
OSPF 改为 RIP & 控制面的信息交换、metric 和收敛行为改变 & 当前 packet 仍使用已安装 FIB\\
HTTP/1.x 改为 HTTP/2/3 & 连接与 stream 的并发组织、队头阻塞边界改变 & 资源/资源表示仍属于应用层\\
\bottomrule
\end{handbooklongtable}

802.1Q 最小帧长度、不同标准/实现的 padding 口径可能不同，涉及数值时必须按题目给定的帧长定义计算，不把“带 VLAN 一定是 68 B”当成无条件规则。

\section{故障题：只追到第一处真正断掉的地方}

看到“网页打不开”“超时”“DNS 成功但 HTTP 没返回”时，不要直接猜自己最熟悉的协议。按照实际过程检查：

\begin{enumerate}
  \item 当前请求是不是已经被本地 HTTP cache 满足？如果没有，主机网络配置是否有效？
  \item 若 URL 使用域名，目标 IP 是否已经得到？
  \item 准备发送的应用数据是否已经形成合法的 TCP/UDP 报文和 IP 数据报？
  \item 当前设备的 FIB 是否给出可执行的出接口和下一跳？
  \item 当前链路需要的下一跳信息是否存在？例如 Ethernet 的 ARP/ND 状态、PPP 的对端状态。
  \item 当前链路帧能否在 VLAN、STP 和 MAC 表的当前状态下完成这一跳？
  \item 经过路由器时，TTL、MTU、NAT 和 LPM 的处理结果是否允许继续转发？
  \item TCP/UDP 以及最终 HTTP 语义是否继续成立？
\end{enumerate}

第一个有证据“不成立”的位置就是\textbf{第一次断点（First Divergence）}。后面的现象可能只是它的连锁结果，不要越过第一处断点继续猜。

\section{贯穿母例：冷启动、热访问和只改一个条件}

\subsection{冷启动：几乎没有可复用状态}

设：DHCP lease 不存在、DNS 缓存未命中、默认网关 ARP 表项不存在、交换机 MAC 表初始为空、TCP 连接不存在、HTTP 缓存未命中，但 STP 和路由已经收敛到可用状态。实际过程可以写成：

\begin{verbatim}
DHCP 建立主机网络配置
-> 得到 IP / 前缀 / 默认网关 / DNS 服务器
-> 需要 DNS 查询
   -> FIB 先选到 DNS server 的下一跳
   -> 需要时 ARP 默认网关
   -> WLAN / Ethernet / routers 把 DNS query 送到 resolver
-> 得到 Web server IP
-> 建立 TCP 连接
   -> 网关 ARP 表项仍有效就直接复用
   -> SYN / SYN+ACK / ACK 每一条都重新走逐跳交付
   -> R1 建立或复用 NAPT 映射
-> HTTP request / response
-> TCP 重组并交付字节
-> HTTP 解释资源表示
-> HTML 还可能引出 CSS / JS / 图片等新请求
\end{verbatim}

最值得注意的不是步骤很多，而是：\textbf{ARP 可以先为了 DNS query 建立，随后继续被 TCP/HTTP 数据复用。}所以 ARP 绝不能固定画在 DNS 之后。

\subsection{热访问：只省掉仍然有效的部分}

第二次很快访问同一站点时，DHCP 租约、DNS 记录、默认网关 ARP 表项、MAC 表、NAT 映射、TCP 连接或 HTTP 缓存可能仍有部分有效。

\begin{itemize}
  \item DNS 缓存命中：不重新发 DNS 查询，但后续 HTTP 数据仍要正常转发；
  \item ARP 缓存命中：不重新发本链路的 ARP Request，但最终 destination IP 不变；
  \item TCP 连接仍可复用：不重新三次握手，但仍要发送需要的 HTTP 请求；
  \item HTTP 缓存仍新鲜并足以满足请求：本次访问可能直接在本地结束，后续网络事件整体消失。
\end{itemize}

\subsection{只改一个条件：只允许相关子过程变化}

这是检验模型是否真正分层的最好方法之一。

\begin{handbooktable}{L{3.2cm} Y Y}
\toprule
只改什么 & 应改变什么 & 不应该跟着改变什么\\
\midrule
DNS 未命中 $\to$ 命中 & 删除 DNS 查询过程 & 后续 Web 数据报仍要查 FIB 并完成链路交付\\
邻居缓存未命中 $\to$ 命中 & 删除当前链路的 ARP/ND 交换 & IP 目的地址、TCP 连接状态\\
TCP 不存在 $\to$ 可复用 & 删除三次握手 & HTTP 请求语义\\
FIB 旧版本 $\to$ 新版本 & 新版本生效后到达的数据报使用新的转发动作 & 已缓存的 DNS 记录\\
MTU 1500 $\to$ 800 & 改变分片、MSS 和每个报文大小相关分支 & 域名本身\\
HTTP 缓存未命中 $\to$ 本地缓存可直接使用 & 后续网络交互可能整体消失 & 后台已经形成的路由/STP 状态\\
\bottomrule
\end{handbooktable}

\section{考纲覆盖检查：没有塞进 Web 主线的知识放在哪里}

主线只选择一组可以同时成立的机制。下面这张表的作用是防止“没有画进主线 = 被删除”的误解。

\begin{handbooklongtable}{L{4.0cm} Y Y}
\toprule
知识点 & 在 NET-I01 中怎样出现 & 完整机制见\\
\midrule
Ethernet / VLAN / 802.11 / PPP framing、FCS、透明传输 & 逐跳路径和字段变化图只保留必要接口 & NET02\\
CSMA/CD、CSMA/CA、polling/token、FDM/TDM/CDM & WLAN 主线 + 替换场景 & NET01 / NET02\\
Switch learning、unknown-unicast flooding、filtering、STP & 接入 LAN 主线 & NET02\\
DHCP DORA/Relay；IPv6 SLAAC/ND & 启动阶段主线 + IPv6 替换场景 & NET04\\
same-subnet 判断、CIDR/LPM、默认路由、ARP & 每条数据报的逐跳交付过程 & NET04；接口见 NET-B01\\
NAPT、TTL、MTU、IPv4 fragmentation、ICMP & R1 边界与异常分支 & NET04；MTU 接口见 NET-B06\\
RIP、OSPF、BGP、收敛中间态 & OSPF/BGP 主场景背景状态；RIP 替换场景 & NET05；接口见 NET-B02\\
端口、UDP header / pseudo-header checksum & DNS/传输层接口 & NET06\\
TCP 三次握手、连接释放、RST、rwnd、零窗口 & 三次握手进入主线；连接关闭在本册停止位置之后 & NET06\\
Stop-and-Wait、GBN、SR、序号空间与窗口 & 作为理解 TCP 可靠传输的前置模型，本册不重复展开 & NET03 / NET-B03\\
slow start、congestion avoidance、fast retransmit/recovery、cwnd & TCP 反馈与可用窗口 & NET07；rwnd/cwnd 接口见 NET-B04\\
BDP 与 window & 性能推导 & NET01 / NET-B05\\
DNS hierarchy、RR、cache TTL & 名字解析主线 & NET08\\
HTTP 资源/资源表示、method/status、cache、对象依赖、HTTP/1.x--3 & 应用阶段与性能题 & NET08\\
FTP control/data connections & Web 主线之外的应用协议替换场景 & NET08\\
SMTP/POP3/IMAP/MIME & 邮件应用协议，不属于 Web 主线 & NET08\\
\bottomrule
\end{handbooklongtable}

\section{做综合题时的最短落笔顺序}

考场不需要复述整本手册。保留下面六步即可：

\begin{enumerate}
  \item \textbf{先写目标和已有状态：}URL/目标是什么？配置、DNS、ARP、FIB、TCP、HTTP cache 哪些已经存在？
  \item \textbf{先判断当前还缺什么：}不要把所有协议机械重跑一遍。
  \item \textbf{一旦要发送一条普通 IP 数据报：}查目的 IP $\to$ 当前 FIB $\to$ 出接口/下一跳 $\to$ 准备本跳链路交付 $\to$ 下一节点。
  \item \textbf{每到一个设备只让它做自己的事：}交换机不减 TTL，路由器不解释 HTTP，DNS 不输出 MAC，TCP 模块不生成路由表。
  \item \textbf{性能题先画依赖关系，故障题先找第一次断点。}
  \item \textbf{最后反查字段：}destination IP、next-hop IP、MAC、port、TTL 有没有被混成同一种“地址”或同一种状态？
\end{enumerate}

\section{独立检查：不用原推理再算一遍}

推演结束后至少做六个反向检查：

\begin{enumerate}
  \item \textbf{名字类型：}domain、destination IP、next-hop IP、MAC、port 有没有被混成同一种“地址”？
  \item \textbf{IP 跳数：}是不是只有真正的 IP 路由器才让 TTL/Hop Limit 推进？
  \item \textbf{链路帧：}到了路由器或明确的链路表示转换边界后，旧二层表示是否还被错误沿用？
  \item \textbf{状态归属：}MAC 表、FIB、ARP cache、TCP state、DNS cache 是否由正确设备或协议模块维护？
  \item \textbf{时间先后：}OSPF/BGP/STP 形成状态、状态生效、当前数据报使用状态，有没有被写成同一时刻？
  \item \textbf{性能代价：}RTT、发送时延和传播时延是否来自真实依赖关系，而不是按协议名机械加常数？
\end{enumerate}

\begin{boundarybox}{本册在哪里停止}
本册到“目标 HTTP 响应的资源表示已经被应用层解释”为止。TCP 四次挥手、半关闭、TIME\_WAIT 和 RST 的完整状态机仍由 NET06 负责；TLS 详细握手、QUIC 内部恢复、CDN 选站、浏览器渲染和内核/NIC 内部执行路径属于其他专题。

停止在这里不是删除这些知识，而是保持职责清楚：NET-I01 负责把一次 Web 请求串起来，不再复制每个局部协议的完整教材。
\end{boundarybox}

\begin{corebox}{最短压缩}
一次网络请求可以记成两句话：

\textbf{先看这次访问还缺什么状态，只补缺失的部分；一旦要发送一条普通 IP 数据报，就用当前设备的 FIB 找出下一跳，并按当前链路把它送过去。}

经过 router 时，入链路帧结束、TTL 推进、下一跳重新决定；经过 AP/switch 时不要误加 IP hop。所有综合题的关键，就是始终分清“最终目的”和“当前下一跳”，分清“哪台设备保存什么状态”，并把“路由/STP 状态怎样形成”与“当前数据报怎样使用这些状态”分开。
\end{corebox}

\end{document}
